\Askhsh{13367}{4}
\begin{ylh}{2}
    \item 4.2. Ανισώσεις 2ου Βαθμού
    \item 6.1. Η Έννοια της Συνάρτησης
    \item 6.2. Γραφική Παράσταση Συνάρτησης
    \item 6.3. Η Συνάρτηση $f(x)=ax+\beta$
\end{ylh}
Δίνεται η ευθεία $\varepsilon$ : $y = \left( \omega^{2} - 6\omega + 8 \right)x + 2$, όπου $\omega\ \varepsilon\ \mathbb{R}$.
\begin{alist}
\item Για τις διάφορες τιμές του $\omega\ \varepsilon\ \mathbb{R}$ να βρείτε το είδος της γωνίας που σχηματίζει η ευθεία $\varepsilon$ με τον άξονα $x'x$.
\monades{9}
\item Αν ο αριθμός $\omega$ είναι ακέραιος και η γωνία που σχηματίζει η ευθεία $\varepsilon$ με τον άξονα $x'x$ είναι αμβλεία τότε:
\begin{rlist}
\item
  Να αποδείξετε ότι $\omega = 3$.
\monades{6}
\item
 Να βρείτε τα σημεία τομής της ευθείας $\varepsilon$ με τους άξονες.

\monades{5}
\item
 Να σχεδιάσετε την ευθεία $\varepsilon$.
\monades{5}
\end{rlist}
\end{alist}

